\documentclass[12pt, a4paper]{article}
\usepackage[T1]{fontenc}
\usepackage[utf8]{inputenc}
\usepackage{amsmath, amssymb, amsthm}
\usepackage{geometry}
\usepackage[hidelinks]{hyperref}
\usepackage{booktabs}
\usepackage{xcolor}
\usepackage{enumitem}
\usepackage{fancyhdr}
\usepackage{pgfplots}
\pgfplotsset{compat=1.18}
\geometry{margin=2.5cm}

\pagestyle{fancy}
\setlength{\headheight}{14pt}
\fancyhf{}
\rhead{Dark Wing -- Financial Mathematics}
\lhead{Rolling Cointegration}
\rfoot{Page \thepage}

\definecolor{darkblue}{RGB}{0,51,102}
\definecolor{alertred}{RGB}{200,0,0}

\title{
    \Huge\textbf{\textcolor{darkblue}{Rolling Window Cointegration}} \\[0.5em]
    \Large Mathematical Notes \& Journal \\[0.5em]
    \large \texttt{rolling\_cointegration.py} — Dark Wing Project
}
\author{Dark Wing — Financial Mathematics Division}
\date{\today}

\begin{document}
\maketitle
\tableofcontents
\newpage

% =============================================================================
\section{The Problem with Static Cointegration}
% =============================================================================

Both the Engle-Granger (\texttt{engle\_granger.py}) and Johansen (\texttt{johansen.py}) tests 
estimate a \emph{single, fixed} cointegrating vector $\hat\beta$ over the entire sample.

This is mathematically convenient but financially dangerous. Markets are non-stationary in 
\emph{structure}, not just in price levels. Consider:

\begin{itemize}
    \item A sector-wide shock (e.g., a budget that hits only one stock)
    \item A regulatory change (e.g., new FPI rules affecting IT but not energy)
    \item Fundamental valuation divergence (e.g., one company grows earnings while the other stagnates)
\end{itemize}

These events can \textbf{break the cointegration relationship permanently or temporarily}. 
A static model will keep signalling trades on a broken pair, leading to sustained losses.

\medskip
\textbf{Rolling cointegration} is the solution: re-test cointegration on a moving window 
of recent data so the model \emph{adapts} to structural changes.

% =============================================================================
\section{Mathematical Setup}
% =============================================================================

\subsection{The Rolling Window}

Let $\{X_t\}_{t=1}^{T}$ be the full dataset of length $T$.

Define a \textbf{window} of size $w$ observations. At each position $n \in \{w, w+1, \ldots, T\}$, 
we extract the sub-sample:

\begin{equation}
    \mathcal{W}_n = \{X_{n-w+1}, X_{n-w+2}, \ldots, X_n\}
    \label{eq:window}
\end{equation}

This is the window \emph{ending} at time $t = n$. It contains the most recent $w$ observations.

For each window $\mathcal{W}_n$, we estimate:
\begin{align}
    \hat\beta_n &= \text{OLS slope of } Y \text{ on } X \text{ using data in } \mathcal{W}_n \label{eq:rolling_beta}\\
    \tau_n      &= \text{ADF } t\text{-statistic on residuals from } \mathcal{W}_n \label{eq:rolling_tau} \\
    p_n         &= \text{p-value of the ADF test at window } n \label{eq:rolling_p}
\end{align}

\subsection{Step Size (Efficiency vs. Precision)}

Running a full cointegration test at \emph{every} time step (step $= 1$) is computationally 
expensive but produces the most granular picture of regime changes.

A \textbf{step} $s > 1$ advances the window by $s$ observations at a time:

\begin{equation}
    n \in \{w,\ w+s,\ w+2s,\ \ldots,\ T\}
    \label{eq:step}
\end{equation}

\textbf{Example:} Daily data, $w = 252$ (1 year), $s = 5$ (weekly step). 
This runs the test once per week, giving $\approx (T - 252)/5$ total tests.

For live production code, step $= 5$ or $= 21$ (weekly / monthly) is a good balance.

% =============================================================================
\section{The Rolling Engle-Granger}
% =============================================================================

At each window $\mathcal{W}_n$, run the two-step Engle-Granger procedure (see \texttt{Engle\_Granger\_Test.tex}):

\begin{align}
    \text{Step 1:}&\quad Y_t = \hat\alpha_n + \hat\beta_n X_t + \hat\varepsilon_t^{(n)} \quad t \in \mathcal{W}_n \label{eq:eg_step1}\\
    \text{Step 2:}&\quad \text{ADF on } \hat\varepsilon_t^{(n)} \rightarrow p_n \label{eq:eg_step2}
\end{align}

The \textbf{output sequence} is a time series of tuples:
\begin{equation}
    \{(\hat\beta_n, p_n, \mathbb{1}_{p_n < \alpha})\}_{n=w}^{T}
\end{equation}

where $\mathbb{1}_{p_n < \alpha}$ is the indicator function (1 if cointegrated, 0 if not) and 
$\alpha$ is the significance level (e.g., 0.05).

% =============================================================================
\section{The Rolling Johansen}
% =============================================================================

For $m \geq 3$ assets, the rolling procedure re-runs the Johansen eigenvalue problem 
(see \texttt{Johansen\_Test.tex}) on each window $\mathcal{W}_n$:

\begin{equation}
    \left| \lambda S_{11}^{(n)} - S_{10}^{(n)} (S_{00}^{(n)})^{-1} S_{01}^{(n)} \right| = 0
\end{equation}

At each window, we record:
\begin{align}
    r_n^* &= \text{cointegrating rank at window } n \\
    \beta_n^{(j)} &= j\text{-th cointegrating vector at window } n
\end{align}

\textbf{Practical note:} For rolling Johansen, we fix $k\_ar\_diff = 1$ inside the rolling 
loop for speed (the optimal lag from the full sample is used as a fixed prior, not recomputed 
at every window).

% =============================================================================
\section{The Stability Score}
% =============================================================================

The \textbf{cointegration stability score} $\mathcal{S}$ quantifies the fraction of the history 
over which the pair was cointegrated:

\begin{equation}
    \boxed{
    \mathcal{S} = \frac{1}{N} \sum_{n=1}^{N} \mathbb{1}_{p_n < \alpha}
    }
    \label{eq:stability}
\end{equation}

where $N$ is the total number of rolling windows computed.

\textbf{Interpretation:}
\begin{center}
\begin{tabular}{@{}cll@{}}
\toprule
\textbf{Stability Score} & \textbf{Classification} & \textbf{Action} \\
\midrule
$\mathcal{S} \geq 0.75$ & Highly Stable  & Good candidate for systematic pairs trading \\
$0.50 \leq \mathcal{S} < 0.75$ & Moderate & Trade with caution; monitor live \\
$\mathcal{S} < 0.50$ & Unstable  & Do not trade algorithmically \\
\bottomrule
\end{tabular}
\end{center}

% =============================================================================
\section{Rolling Hedge Ratio: Time-Varying $\hat\beta_n$}
% =============================================================================

The sequence $\{\hat\beta_n\}_{n=w}^{T}$ is the \textbf{rolling hedge ratio}. 
It tells you how the long-run equilibrium relationship has evolved.

\subsection{Why This Matters for Position Sizing}

When you construct the spread $S_t = Y_t - \hat\beta_n \cdot X_t$, you should use the 
\emph{most recent} estimate of $\hat\beta_n$ rather than the static full-sample $\hat\beta$.

\medskip
Using an outdated $\hat\beta$ creates a \textbf{stale hedge}: the portfolio is no longer 
market-neutral, and the "spread" you are trading will have a directional bias.

\subsection{Relationship to the Kalman Filter}

The rolling hedge ratio $\hat\beta_n$ is a \emph{batch} estimate — it re-estimates the hedge 
ratio from scratch using the last $w$ data points.

The \textbf{Kalman filter} (\texttt{dynamic\_beta.py}) is the online, sequential equivalent: 
it updates $\hat\beta_t$ one data point at a time using a state-space model. 

\begin{center}
\begin{tabular}{@{}lll@{}}
\toprule
\textbf{Method} & \textbf{Approach} & \textbf{Best for} \\
\midrule
Rolling EG (this file)  & Batch re-estimation   & Regime detection, strategy research \\
Kalman Filter           & Sequential update     & Live trading, real-time signal \\
\bottomrule
\end{tabular}
\end{center}

In the \textbf{Dark Wing pipeline}, both are used together:
\begin{enumerate}
    \item \textbf{Rolling cointegration} determines if a pair is worthy of trading (stability score).
    \item \textbf{Kalman filter} provides the real-time hedge ratio for actual trade execution.
\end{enumerate}

% =============================================================================
\section{Structural Breaks and Regime Classification}
% =============================================================================

A structural break in cointegration occurs when $p_n$ crosses the significance threshold $\alpha$:

\begin{equation}
    \text{Break at time } n \iff p_{n-1} < \alpha \text{ and } p_n \geq \alpha
    \label{eq:break}
\end{equation}

A \textbf{recovery} (re-establishment of cointegration) occurs when:
\begin{equation}
    \text{Recovery at time } n \iff p_{n-1} \geq \alpha \text{ and } p_n < \alpha
    \label{eq:recovery}
\end{equation}

\textbf{In practice:} Do not immediately trade a break or recovery signal. 
Apply a \textbf{cooldown rule}: the pair must be continuously cointegrated for at least 
$c$ consecutive windows (e.g., $c = 10$) before re-entering a position.

% =============================================================================
\section{Window Size Selection}
% =============================================================================

The choice of window size $w$ involves a fundamental bias-variance tradeoff:

\begin{center}
\begin{tabular}{@{}cll@{}}
\toprule
\textbf{Window Size} & \textbf{Advantage} & \textbf{Disadvantage} \\
\midrule
Small ($w < 100$)  & Adapts quickly to regime changes & High variance; unstable $\hat\beta$ \\
Large ($w > 500$)  & Stable, reliable estimates        & Slow to detect breaks \\
Medium ($w \approx 252$) & Balance of both              & Most practical for daily data \\
\bottomrule
\end{tabular}
\end{center}

\textbf{Rule of thumb for stocks:} $w = 252$ (1 trading year) is the standard for daily data.
For intraday data, adjust proportionally (e.g., $w = 252 \times 78$ for 5-minute bars on 252 days).

% =============================================================================
\section{Connection to the Dark Wing Pipeline}
% =============================================================================

\begin{enumerate}[label=\textbf{Step \arabic*:}]
    \item \textbf{Pair Screening:} Run Engle-Granger or Johansen on all candidate pairs 
          using the full historical dataset.
    \item \textbf{Stability Filter:} Run \texttt{RollingCointegration.run()} and compute 
          the stability score $\mathcal{S}$ for each pair.
    \item \textbf{Shortlist:} Only pairs with $\mathcal{S} \geq 0.75$ are approved for live trading.
    \item \textbf{Current Status:} Use \texttt{RollingCointegration.current\_status()} to check 
          if the pair is cointegrated \emph{right now} before placing any order.
    \item \textbf{Live Execution:} For shortlisted, currently-cointegrated pairs, launch the 
          Kalman filter (\texttt{dynamic\_beta.py}) to track $\beta_t$ in real-time.
    \item \textbf{Monitoring:} Re-run the rolling cointegration weekly (via a scheduler) to 
          detect regime breaks and halt trading automatically.
\end{enumerate}

% =============================================================================
\section{Summary of Key Formulas}
% =============================================================================

\begin{align}
    \text{Window at } n:& \quad \mathcal{W}_n = \{X_{n-w+1}, \ldots, X_n\} \tag{\ref{eq:window}} \\[6pt]
    \text{Rolling OLS:}& \quad \hat\beta_n = \frac{\sum_{t \in \mathcal{W}_n}(Y_t - \bar Y_n)(X_t - \bar X_n)}{\sum_{t \in \mathcal{W}_n}(X_t - \bar X_n)^2} \tag{\ref{eq:rolling_beta}} \\[6pt]
    \text{ADF on residuals:}& \quad \Delta\hat\varepsilon_t = \gamma \hat\varepsilon_{t-1} + \sum_j c_j \Delta\hat\varepsilon_{t-j} + u_t \tag{\ref{eq:eg_step2}} \\[6pt]
    \text{Stability Score:}& \quad \mathcal{S} = \frac{1}{N}\sum_{n=1}^N \mathbb{1}_{p_n < \alpha} \tag{\ref{eq:stability}} \\[6pt]
    \text{Structural Break:}& \quad p_{n-1} < \alpha \text{ and } p_n \geq \alpha \tag{\ref{eq:break}}
\end{align}

\end{document}
